\chapter{Introducción}

\section{Objetivos generales}

Este Trabajo de Fin de Grado se enmarca en el desarrollo de un laboratorio
virtual para la simulación y el análisis de paradigmas de toma de decisiones
humanas mediante agentes inteligentes. El proyecto completo se divide en dos
fases complementarias. Esta memoria documenta la fase desarrollada por Pablo
Pazos Parada: el sistema que genera modelos de decisión ejecutables a partir de
descripciones en lenguaje natural y conocimiento científico.

El objetivo general de esta fase es construir un pipeline capaz de transformar
un problema de decisión en un conjunto de artefactos verificables:

\begin{itemize}
\item una investigación sobre paradigmas científicos relevantes;
\item una o varias formulaciones matemáticas;
\item una especificación adaptada a un entorno de simulación;
\item código Python autocontenido;
\item pruebas ejecutables que validen el modelo generado.
\end{itemize}

\section{Relación de la documentación}

La memoria sigue la estructura de un TFG de tipo B, orientado al desarrollo de
un sistema software. Los capítulos principales son:

\begin{itemize}
\item especificación de requisitos;
\item diseño del sistema;
\item desarrollo del trabajo;
\item pruebas;
\item conclusiones y posibles ampliaciones.
\end{itemize}

Los apéndices incluyen información técnica de ejecución y una guía de uso
básica. La bibliografía recoge tanto referencias técnicas de agentes y memoria
como trabajos de recuperación aumentada, grafos de conocimiento y generación de
código.

\section{Descripción del sistema}

El sistema recibe una descripción de un problema de toma de decisiones y la
procesa mediante agentes especializados. Cada agente resuelve una parte del
problema y produce artefactos persistentes que pueden ser revisados antes de
continuar.

\begin{figure}[H]
\centering
\begin{verbatim}
Descripcion del problema
        |
        v
Investigacion cientifica
        |
        v
Formalizacion matematica
        |
        v
Adaptacion al entorno
        |
        v
Modelo Python + pruebas
\end{verbatim}
\caption{Transformación principal del sistema desarrollado.}
\end{figure}

La salida final es un modelo compatible con el contrato de simulación usado por
el laboratorio virtual:

\begin{lstlisting}[language=Python]
def decide(self, perception: dict) -> Action
def update(self, action, reward, new_perception) -> None
def get_state(self) -> dict
\end{lstlisting}

Esta interfaz permite separar la fase de generación de modelos de la fase de
simulación. El modelo generado no necesita heredar de una clase concreta del
simulador; basta con que implemente los métodos esperados.

\section{Información adicional de interés}

El trabajo combina varias técnicas actuales:

\begin{itemize}
\item agentes de lenguaje con uso de herramientas;
\item revisión humana entre etapas;
\item almacenamiento de artefactos;
\item recuperación aumentada con búsqueda densa, BM25 y grafo;
\item memoria persistente con ciclo de vida;
\item generación de código con pruebas automáticas.
\end{itemize}

Una característica importante del proyecto es que el propio desarrollo del
sistema también se realizó con apoyo intensivo de agentes de programación. Esta
metodología se describe de forma específica en el capítulo de desarrollo, porque
afecta a la planificación, a la división de tareas, a la validación y a la forma
de tomar decisiones de diseño.
